%%=============================================================================
%% Voorwoord
%%=============================================================================

\chapter*{\IfLanguageName{dutch}{Woord vooraf}{Preface}}%
\label{ch:voorwoord}

%% TODO:
%% Het voorwoord is het enige deel van de bachelorproef waar je vanuit je
%% eigen standpunt (``ik-vorm'') mag schrijven. Je kan hier bv. motiveren
%% waarom jij het onderwerp wil bespreken.
%% Vergeet ook niet te bedanken wie je geholpen/gesteund/... heeft

In deze bachelorproef onderzoek ik hoe apps voor stress- en burn-outpreventie het mentale welzijn van Vlaamse twintigers kunnen ondersteunen. Dit onderzoek vormt het sluitstuk van mijn opleiding toegepaste informatica, waarna ik hoop af te studeren.

\vspace{1em}

De keuze voor dit onderwerp kwam niet zomaar tot stand. Tijdens de COVID-19-pandemie, gedurende mijn vorige opleiding, merkte ik hoe een hoge werkdruk en te weinig rust hun tol konden eisen. Ik ging daarom zelf op zoek naar digitale toepassingen die konden helpen om meer rust te vinden. Helaas boden deze apps mij niet de ondersteuning die ik ervan verwachtte, waardoor ik ze al snel weer links liet liggen. Deze ervaring deed mij nadenken over hoe zulke toepassingen beter kunnen worden ontworpen. Tijdens deze bachelor ontdekte ik dat mentale gezondheidsapps sterk zijn gegroeid sinds de pandemie, terwijl stress- en burn-outklachten bij jongvolwassenen bleven toenemen. Vanuit die vaststelling ontstond het idee voor deze bachelorproef: onderzoeken waar bestaande toepassingen tekortschieten en deze inzichten vertalen naar richtlijnen voor meer stressreducerende mentale gezondheidsapps.

\vspace{1em}

Het uitvoeren van dit onderzoek was intensief maar bijzonder leerrijk. De combinatie van literatuuronderzoek, requirementsanalyse, mock-ups en de implementatie van de \acrshort{poc} bracht verschillende uitdagingen met zich mee. Niet alle ideeën konden binnen de beschikbare tijd volledig worden uitgewerkt. De literatuurstudie deed mij bovendien inzien hoe belangrijk psychologische basisbehoeften zoals autonomie, competentie en verbondenheid zijn bij het ontwerpen van digitale toepassingen. Dit verbreedde mijn kijk op softwareontwikkeling en maakte de psychologische invalshoek van het onderzoek voor mij bijzonder waardevol. Dankzij mijn copromotor kon ik hierbij rekenen op waardevolle begeleiding en inzichten vanuit de psychologie.

\vspace{1em}

Daarvoor wil ik mijn copromotor, mevrouw K. Van den Broeck, graag bedanken. Ook mijn promotor, mevrouw C. De Leenheer, wil ik bedanken voor haar feedback en begeleiding gedurende het volledige traject. Daarnaast bedank ik meneer T. Parmentier om tijdelijk als vervangende promotor in te vallen. Ten slotte wil ik mijn ouders, mijn zus A. Emanuel en mijn vriend M. Dhondt bedanken voor hun voortdurende steun, motivatie en aanmoediging tijdens het schrijven van deze bachelorproef. Dankzij hen stond ik er tijdens dit traject nooit alleen voor.